\subsection{6. Implementation}\label{implementation}

This section presents the implementation using the Stable Baselines3
framework in the \texttt{experiments-codesource/} folder.

\hypertarget{data-preprocessing-helper_env.py}{%
	\subsection{6.1 Data Preprocessing (helper\_env.py)}\label{data-preprocessing-helper_env.py}}

The \texttt{STRCalculator} class implements hierarchical distance → STR
→ capacity calculation:

\textbf{STRCalculator}: Implements Signal Transmission Reward
calculation based on Euclidean distance with exponential attenuation: -
\texttt{gps\_to\_enu()}: Convert GPS to East-North-Up coordinates -
\texttt{enu\_to\_polar()}: Convert to ego-centric polar coordinates -
\texttt{compute\_capacity()}: Shannon-Hartley capacity based on STR
model - \texttt{compute\_rewards()}: Reward calculation from capacity -
\texttt{get\_reshaped\_states()} / \texttt{get\_reshaped\_rewards()}:
State and reward preprocessing - \texttt{fill\_states\_columns()}: State
padding for variable AP counts

\hypertarget{simulation-environment-illinois_online.py}{%
	\subsubsection{6.2 Simulation Environment
		(illinois\_online.py)}\label{simulation-environment-illinois_online.py}}

\textbf{MovingIoTEnvironment} (\texttt{illinois\_online.py}):
Gymnasium-compliant simulation environment for moving IoT service
composition with service mobility following random waypoint model,
device mobility with boundary reflection, and STR-based reward
calculation. Supports both online and offline modes.

\hypertarget{stable-baselines3-framework-for-a2c-and-dqn-training}{%
	#### 6.3 Stable Baselines3 Framework for A2C and DQN Training
}\label{stable-baselines3-framework-for-a2c-and-dqn-training}

We utilize the Stable Baselines3 (SB3) framework
(\url{https://stable-baselines3.readthedocs.io/}) for both A2C and DQN
implementations, replacing custom PyTorch implementations with
well-tested, production-ready algorithms.

\textbf{Installation}:

\begin{verbatim}
pip install stable-baselines3
\end{verbatim}

\textbf{A2C Implementation (\texttt{train.py})}: We use
\texttt{stable\_baselines3.A2C} with the custom Gymnasium environment:

\begin{figure*}[t]
	\centering
	\caption{A2C Implementation Code}
	\label{a2c-implementation-code}
	\scriptsize
	\begin{verbatim}
from stable_baselines3 import A2C
from illinois_online import APSelectionEnv

env = APSelectionEnv(num_services=20, max_steps=100)

model = A2C(
    policy="MlpPolicy",
    env=env,
    learning_rate=5e-5,
    n_steps=30,
    gamma=0.9,
    ent_coef=0.05,
    vf_coef=0.5,
    max_grad_norm=1.0,
    verbose=1
)

model.learn(total_timesteps=100000)
model.save("a2c_moving_iot")
	\end{verbatim}
\end{figure*}

\textbf{DQN Implementation}: We use \texttt{stable\_baselines3.DQN}:

\begin{verbatim}
from stable_baselines3 import DQN

model = DQN(
    policy="MlpPolicy",
    env=env,
    learning_rate=5e-5,
    buffer_size=100000,
    learning_starts=1000,
    gamma=0.9,
    target_update_interval=500,
    exploration_fraction=0.1,
    verbose=1
)

model.learn(total_timesteps=100000)
model.save("dqn_moving_iot")
\end{verbatim}

\textbf{Configuration Management (\texttt{config\_mgmt.py})}: -
\texttt{MasterConfig}: Pydantic model for all hyperparameters -
\texttt{load\_config()}: YAML-based config loading - Supports train/eval
phases with different num\_services settings

\textbf{Hyperparameters} (from YAML config):

\begin{tabular}{|l|l|}
	\hline
	\textbf{Parameter}     & \textbf{Value} \\
	\hline
	Learning Rate          & 5e-5           \\
	\hline
	Discount Factor ($\gamma$)    & 0.9            \\
	\hline
	N-step (A2C)           & 30             \\
	\hline
	Entropy Coefficient    & 0.05           \\
	\hline
	Buffer Size (DQN)      & 100000         \\
	\hline
	Target Update Interval & 500            \\
	\hline
	Max Grad Norm          & 1.0            \\
	\hline
	Episodes               & 100            \\
	\hline
	Target Accuracy        & 98\%           \\
	\hline
\end{tabular}

The real GPS experiments use a production-ready automation system for
reproducible research:

\textbf{Workflow}:
1. \textbf{Configuration} (\texttt{configs/*.yaml}): Define experiment parameters
2. \textbf{Execution} (\texttt{train.py}): Load config, setup directories, run experiment
3. \textbf{Tracking} (\texttt{utils.py}): Log metrics, save results to \texttt{experiments\_results.csv}
4. \textbf{Visualization} (\texttt{training\_plots.py}): Generate training curves and action distributions

\textbf{Key Scripts}:

\begin{verbatim}
# Single experiment
python train.py --config configs/exp1.yaml

# Multiple experiments (background)
./run_experiments.sh configs/exp1.yaml configs/exp2.yaml
\end{verbatim}

\textbf{Output Structure}:

\begin{verbatim}
runs/<run_id>/
  +-- model/
  |   +-- a2c_model.zip          # SB3 A2C model checkpoint
  |   +-- resolved_config.yaml   # Resolved configuration
  +-- plot/
  |   +-- episode_rewards.png    # Training metrics plot
  |   +-- action_distribution.png # Action distribution
  +-- logs/
      +-- <run_id>.log           # Execution logs

experiments_results.csv        # Aggregated results
\end{verbatim}

\subsection{7. Discussion}\label{discussion}

\hypertarget{interpretation-of-results}{%
	\subsubsection{7.1 Interpretation of Results}\label{interpretation-of-results}}

A2C outperforms Double DQN on moving IoT service composition while building upon the STR-based selection from prior work \cite{neiat2021}. Three factors drive improvements.

First, actor-critic gives stable learning by combining value estimation
with direct policy optimization. The advantage function reduces variance
while keeping updates unbiased, so the agent learns effectively from
fewer samples.

Second, the relative polar coordinate representation captures spatial
relationships efficiently, enabling the agent to learn effective service
selection without explicit velocity or trajectory prediction.

Third, separate networks specialize processing for spatio-temporal
states. More parameters and training time, but better performance when
movement patterns matter.

The STR-based selection continues driving decisions---the agent learns
to maximize expected capacity while maintaining stability.

\hypertarget{comparison-with-prior-work}{%
	\subsubsection{7.2 Comparison with Prior
		Work}\label{comparison-with-prior-work}}

A2C improves over Double DQN on both datasets.

ATC (Indoor): A2C Separate achieves 95.2\% versus 92.4\% at low
mobility (2.8\% improvement). At high mobility (10 km/h), the gap grows
to 13.9\% (85.7\% vs 71.8\%).

Illinois (Outdoor): At high mobility (10 km/h), A2C Separate reaches
80.3\% versus 64.2\% for Double DQN---a 16.1\% absolute improvement.
Outdoor trajectories with more varied movement patterns benefit from
the learned policy.

Capacity satisfaction improves because A2C better captures spatial
relationships and selects services with better capacity margins.

\hypertarget{practical-implications}{%
	#### 7.3 Practical Implications}\label{practical-implications}}

\textbf{Edge Deployment}: A2C with shared networks balances performance
and computation for edge deployment. Inference cost allows real-time
decisions on edge devices.

\textbf{Stability}: A2C re-composes less frequently than Double DQN,
reducing service disruption and overhead for production systems.

\textbf{STR Quality}: STR-derived capacity provides a grounded service
quality metric tied to spatial proximity.

\hypertarget{limitations}
#### 7.4 Limitations}\label{limitations}

Three limitations point to future work:

\textbf{STR Model}: Uses simplified distance-based models. More complex
propagation---multipath fading, shadowing, interference---needs
investigation.

\textbf{Centralized}: Assumes centralized composition. Distributed
multi-agent approaches could scale better for large IoT systems.

\textbf{Validation}: We tested on real GPS from ATC and Illinois
datasets, but service availability is simulated. Future work should
use real IoT service availability data.

\begin{center}\rule{0.5\linewidth}{0.5pt}\end{center}

\hypertarget{conclusion}{%

